\begin{frame}
\frametitle{Die Beschreibung von Verteilungen}

Nachdem wir operationalisiert und Daten für die theoretischen Konstrukte erhoben haben, liegen die Daten nun in einer Datenmatrix vor.\footnote{Aufgebaut wie eine Excel-Tabelle: Die Zeilen sind die Untersuchungseinheiten und die Spalten sind die Eigenschaften. Jede Zelle beinhaltet daher die Ausprägung einer Eigenschaft für eine Untersuchungseinheit.} Diese Datenmatrix ist meist in ihrer Gesamtheit uninformativ, denn einzelne Werte sagen uns nichts über die \textit{Typik} von Eigenschaften oder über Zusammenhänge zwischen Eigenschaften etwas aus. \vfill

Um Daten gut zu beschreiben, müssen wir sie sinnvoll zusammenfassen (= aggregieren). Dies ist Aufgabe der \textit{deskriptiven Statistik}. 


\end{frame}
\begin{frame}
\frametitle{Was ist eine Verteilung?}

Eine Verteilung der Werte einer Eigenschaft $Y$ gibt an, wie \textbf{häufig} diese Werte innerhalb einer (analytischen) Stichprobe vorkommen. Diese Häufigkeit kann in unterschiedlicher Weise angegeben werden: 

\vfill

\begin{enumerate}
  \item absolute Häufigkeiten
  \item relative Häufigkeiten / Wahrscheinlichkeit 
  \item Dichte
\end{enumerate}

\vfill

Diese Häufigkeiten werden auf der Y-Achse (der vertikalen Dimension) grafisch abgetragen. Auf der X-Achse (horizontale Dimension) werden die Ausprägungen abgetragen. Je nach Skalenniveau der gemessenen Eigenschaften sind dies daher Kategorien, Ränge oder Werte in gleichen Intervallen. 

\end{frame}
\begin{frame}
\frametitle{Darstellung von Verteilungen als Balkendiagramme}
\begin{columns}
\begin{column}{.6\linewidth}
\resizebox{1\linewidth}{!}{\includegraphics{lehre-gew-arbeitsstunden-cat}}
{{ \scriptsize Verteilung der gewünschten Arbeitsstunden, 2018. Daten: SOEP V35.}}
\end{column}
\begin{column}{.4\linewidth}
Verteilungen mit wenigen Ausprägungen können in Balkendiagrammen dargestellt werden.\footnote{Ich sage es nur einmal: Wer Tortendiagramme oder 3D-Diagramme anfertigt, kommt in die Hölle. }\newline 

\vfill 

Achten Sie bei Balkendiagrammen auf eine sinnvolle Abtragung der Y-Achse. Diese muss nicht zwingend vom Minimalwert bis zum Maximalwert gehen.\newline   

\end{column}
\end{columns}
\end{frame}
\begin{frame}
\frametitle{Darstellung von Verteilungen als Balkendiagramme}
\begin{columns}
\begin{column}{.6\linewidth}
\resizebox{1\linewidth}{!}{\includegraphics{lehre-gew-arbeitsstunden-dis}}
{{ \scriptsize Verteilung der gewünschten Arbeitsstunden, 2018. Daten: SOEP V35.}}
\end{column}
\begin{column}{.4\linewidth}
Je mehr Ausprägungen der Eigenschaft Sie darstellen wollen, desto unübersichtlicher werden Balkendiagramme für die Darstellung von Verteilungen. Das ist insbesondere für metrisch gemessene Eigenschaften der Fall, die viele (oder potentiell: unendlich viele) unterschiedliche Ausprägungen haben können.  
\end{column}
\end{columns}
\end{frame}
\begin{frame}
\frametitle{Dastellung von Verteilungen als Kurven}
\begin{columns}
\begin{column}{.5\linewidth}
\resizebox{1\linewidth}{!}{\includegraphics{lehre-gew-arbeitsstunden-kdens}}
{{ \scriptsize Verteilung der gewünschten Arbeitsstunden, 2018. Quelle: SOEP V35.}}
\end{column}
\begin{column}{.5\linewidth}
\small 
Metrische Verteilungen können auch als Kurven oder Funktionen abgetragen werden. Dabei werden immer (!) die zugrunde liegenden Daten geglättet (= Smoothing). Je nach Stärke der Glättung kann sich die Darstellung der Verteilung stark ändern. Je stärker geglättet wird, desto mehr werden die zugrunde liegenden Daten zu einer (meist glockenförmigen) Verteilung angenähert.\\ \newline

\vfill 

Sehr glatte und glockenförmige Verteilungen sind empirisch kaum beobachtbar. Empirische Verteilungen haben typischerweise ``Zacken'' oder ``Klumpen''. Das ist kein Problem, sondern Teil des empirischen Musters, dass Sie studieren wollen. 

\end{column}
\end{columns}
\end{frame}
\begin{frame}
\frametitle{Die Beschreibung von Verteilungen}
\begin{columns}
\begin{column}{.5\linewidth}
\resizebox{1\linewidth}{!}{\includegraphics{lehre-gew-arbeitsstunden-kdens-gender}}
{{ \scriptsize Verteilung der gewünschten Arbeitsstunden nach Geschlecht, 2018. Daten: SOEP V35.}}
\end{column}
\begin{column}{.5\linewidth}
\small 
Verteilungen können sich in mehreren Aspekten unterscheiden: 
\begin{enumerate}
  \item Werte können häufiger sein
  \item der häufigste Wert kann sich unterscheiden
  \item der ``zentrale'' Wert kann sich unterscheiden
  \item die Werte können unterschiedlich breit streuen
  \item die Werte können symmetrisch um einen zentralen Wert streuen oder wenige Werte können über einen sehr großen Bereich streuen. 
\end{enumerate}

Unser Ziel ist es, diese Unterschiede in weniger, klar verständlichen Kennzahlen auszudrücken. 

\end{column}
\end{columns}
\end{frame}
\section{Häufigkeiten}
\subsection*{}
\begin{frame}
  \frametitle{Häufigkeiten: Notation \& Begriffe}
\small
\begin{tabular}{llp{.5\linewidth}} \toprule 
Notation & Bedeutung & Beispiel\\ \midrule
$N$ & Absolute Häufigkeit & 4.850 Personen gaben in der Stichprobe von 2018 an, 40h arbeiten zu wollen. Die Stichprobe besteht aus  16.886 Personen. \\[.5em] 
$P(y)$ & Wahrscheinlichkeit & Mit einer Wahrscheinlichkeit von 0.2872 ziehen wir aus der Stichprobe eine Person, die genau 40h als gewünschte Arbeitszeit angegeben hat.\\[.5em] 
$P(y)\cdot 100$ & relative Häufigkeit & 28.72\% der Befragten gaben an, 40h arbeiten zu wollen.\\[.5em]  
$P(y) \, | \, (X = x)$ & Bedingte Wahrscheinlichkeit & Unter der Bedingung, dass wir nur Männer analysieren, ist die Wahrscheinlichkeit, einen Befragten zu ziehen, der 40h arbeiten will, 0.4. Bei Frauen liegt diese Wahrscheinlichkeit bei 0.17\\ \bottomrule
\end{tabular}
  
\end{frame}
\frame{
\frametitle{Wahrscheinlichkeit \& relative Häufigkeit}

Die Wahrscheinlichkeit $P(y)$ einer Ausprägung eines Merkmals Y in Bezug auf alle Beobachtungen N lässt sich daher wie folgt berechnen: 

\begin{equation}  
  P(y) = \frac{y}{N}
\end{equation}

\vfill

Eine Wahrscheinlichkeit von 0,2 sagt Ihnen, dass 20 von 100 Untersuchungseinheiten eine Ausprägung $\tau$ der Eigenschaft Y aufweisen.  Relative Häufigkeiten sind Wahrscheinlichkeiten multipliziert mit 100. 

}
\frame{
\frametitle{Bedingte Wahrscheinlichkeit \& relative Häufigkeit}

Bedingte Wahrscheinlichkeiten oder bedingte relative Häufigkeiten werden innerhalb der durch eine Bedingung festgelegten Gruppen gebildet. 

\begin{equation}  
  P(y) \, | (X = x) = \frac{y}{N_x}
\end{equation}

\vfill

Ist X zum Beispiel eine Variable für Bildung mit drei Ausprägungen (1 = ohne Berufsabschluss, 2 = mit Berufsabschluss, 3 = Hochschulabschluss) und y sind die Angaben zu den gewünschten Arbeitsstunden, geben die bedingten Wahrscheinlichkeiten an, wie wahrscheinlich bestimmte Arbeitsstunden innerhalb der jeweiligen Bildungsgruppe gefunden werden können. 

}
\begin{frame}
\frametitle{Relative Häufigkeiten drücken Relationen aus!}

Wahrscheinlichkeiten oder \textit{relative} Häufigkeiten haben immer eine Bedeutung in Relation zum Nenner. Wenn sich der Nenner ändert, ändert sich auch die Bedeutung der relativen Häufigkeit. Klären Sie daher stets, auf welche \textit{Basis} die relativen Häufigkeiten beziehen. Wir können folgende Relationen grob unterscheiden: 

\begin{itemize}
   \item Relation zur Stichprobengröße 
   \item Relation zu einer bestimmten Gruppe definiert durch Ausprägungen in X
   \item Relation zu einer bestimmten Gruppe definiert durch Ausprägungen in Y
 \end{itemize} 

Die erste Variante wird typischerweise als ``Zellenprozent'', die zweite als ``Zeilenprozent'' und die dritte als ``Spaltenprozent'' bezeichnet. 

\end{frame}
\begin{frame}[fragile]
\frametitle{Armutsrisiko (Zeilenprozent)}
\begin{columns}
\begin{column}{.6\linewidth}
\begin{Verbatim}[fontsize=\small]
Erwerbskonstella |  Armutsgef. Haushalt
            tion |      nein         ja |     Total
-----------------+----------------------+----------
         Rentner |     84.74      15.26 |    100.00 
      Joblose HH |     33.51      66.49 |    100.00 
 Doppelverdiener |     99.61       0.39 |    100.00 
 Alleinverdiener |     91.42       8.58 |    100.00 
     Zuverdiener |     97.68       2.32 |    100.00 
  Doppelteilzeit |     88.03      11.97 |    100.00 
      Prekäre HH |     52.88      47.12 |    100.00 
Single, Vollzeit |     92.24       7.76 |    100.00 
Single, Teilzeit |     62.32      37.68 |    100.00 
-----------------+----------------------+----------
           Total |     83.14      16.86 |    100.00 

\end{Verbatim}
\end{column}
\begin{column}{.3\linewidth}
Wie hoch ist die Wahrscheinlichkeit einen Haushalt in Armut zu finden \textit{innerhalb} eines bestimmten Haushaltstyps? 
\end{column}
\end{columns}

\end{frame}
\begin{frame}[fragile]
\frametitle{Zusammensetzung von Armut (Spaltenprozent)}
\begin{columns}
\begin{column}{.6\linewidth}
\begin{Verbatim}[fontsize=\small]
Erwerbskonstella |  Armutsgef. Haushalt
            tion |      nein         ja |     Total
-----------------+----------------------+----------
         Rentner |     37.82      33.58 |     37.10 
      Joblose HH |      2.79      27.30 |      6.92 
 Doppelverdiener |     10.24       0.20 |      8.55 
 Alleinverdiener |      7.10       3.28 |      6.46 
     Zuverdiener |     12.63       1.48 |     10.75 
  Doppelteilzeit |      1.10       0.74 |      1.04 
      Prekäre HH |      2.03       8.93 |      3.20 
Single, Vollzeit |     21.00       8.72 |     18.93 
Single, Teilzeit |      5.29      15.78 |      7.06 
-----------------+----------------------+----------
           Total |    100.00     100.00 |    100.00 
\end{Verbatim}
\end{column}
\begin{column}{.3\linewidth}
Wie hoch ist die Wahrscheinlichkeit einen bestimmten Haushaltstyp zu finden \textit{innerhalb} der armutsbetroffenen Haushalte? 
\end{column}
\end{columns}

\end{frame}

\section{Typische Werte}

\begin{frame}
  \frametitle{Typische Werte: Notation \& Begriffe}
\small
\begin{tabular}{lp{.12\linewidth}p{.3\linewidth}p{.4\linewidth}} \toprule 
Notation & Bezeichnung & Bedeutung & Beispiel\\ \midrule
$y_{Mod} $ & Modalwert & Der Wert mit der höchsten Häufigkeit & Die häufigsten gewünschten Arbeitsstunden waren 40h. \\[.5em]
$\tilde{y}$ & Median & Der Wert der, die Verteilung in zwei gleich große relative Häufigkeiten teilt & 50\% der Befragten wünschten sich eine Arbeitszeit von unter 35h und 50\% wünschten sich eine Arbeitszeit ab 35h.\\[.5em]
$\bar{y}$ & Arith\-metisches Mittel & Der Wert, der die Verteilung in zwei gleich große Wertesummen teilt. & Das arith. Mittel der gewünschten Arbeitsstunden lag im Jahr 2018 bei 32.8h. Die Summe aller gewünschten Arbeitsstunden unter 32.8h ist genauso groß wie die Summe aller Arbeitsstunden darüber.\\
 \bottomrule
\end{tabular}
  
\end{frame}
\begin{frame}
\frametitle{Modalwert}
\begin{columns}
\begin{column}{.5\linewidth}
\begin{center}
\begin{tabular}{c} \\ \toprule
Größe in cm\\ \midrule
160\\
\textbf{165}\\
\textbf{165}\\
\textbf{165}\\
 173\\ 
 174\\
 176\\
 180\\
 180\\
 187\\ \bottomrule
\end{tabular}
\end{center}
\quad \newline 
\end{column}
\begin{column}[b]{.5\linewidth}
Der Modalwert oder Modus ist der \textbf{häufigste} Wert einer Verteilung. Es ist daher der Wert mit der größten relativen Häufigkeit. 

\[
y_{Mod} = 165  \nonumber 
\]
\end{column}
\end{columns}
\quad \newline 
\end{frame}
\begin{frame}
\frametitle{Median}
Der Median ist der Wert der Beobachtung einer Verteilung, der genau in ihrer Mitte liegt. 

Der Median $\tilde{y}$ von n \textbf{nach der Größe geordneten} Beobachtungen berechnet sich durch: 

\begin{equation}
  \tilde{y} = \left\{  
\begin{array}{l l}
 y_{\frac{n+1}{2}} \quad & \text{bei ungerader Anzahl}\\[1em]
 \frac{y_\frac{n}{2} + y_{\frac{n}{2}+1}}{2} \quad & \text{bei gerader Anzahl}\\
 \end{array} \right.
\end{equation}

$n$ bezeichnet die Anzahl der Beobachtungen. Bei 10 Beobachtungen bezeichnet $y_\frac{10}{2} = y_5 $ den \textbf{Wert} der fünften Beobachtung innerhalb der nach den \textit{Y-Werten geordneten Verteilung}. 

\end{frame}
\frame{
\frametitle{Median}

\begin{center}
  \begin{minipage}[b]{.4\linewidth} % [b] => Ausrichtung an \caption
  
\begin{center}
\begin{tabular} {@{}l c @{}} \\ \hline
Nr. & Größe in cm\\
\hline
1 &160\\
2 &165\\
3 &165\\
4 & 165\\
\textbf{5} & 173\\ \hline
\textbf{6} & 174\\
7 & 176\\
8 & 180\\
9 & 180\\
10 &187\\
\hline
\end{tabular}
\vspace{1em}
\end{center}
  \end{minipage}
  %\hspace{.1\linewidth}% Abstand zwischen Bilder
  \begin{minipage}[b]{.4\linewidth} % [b] => Ausrichtung an \caption
\begin{align}
 \tilde{y} &= \frac{y_\frac{n}{2} + y_{\frac{n}{2}+1}}{2} \nonumber \\  
           &= \frac{y_\frac{10}{2} + y_{\frac{10}{2}+1}}{2} \nonumber \\
           &= \frac{y_{5} + y_{6}}{2} \nonumber \\
           &= \frac{173 + 174}{2} = 173,5 \nonumber
 \end{align}
  \end{minipage}
\end{center}

}
\frame{
\frametitle{Arithmetisches Mittel}

Das arithmetische Mittel $\overline{y}$ der Werte $y_i$ berechnet sich in Bezug auf die Größe einer Verteilung $N$ grundsätzlich in folgender Art: 

\begin{equation} 
\bar{y} = \dfrac{\sum y_i}{N}
\end{equation}

}
\frame{
\frametitle{Arithmetisches Mittel}

\begin{center}
  \begin{minipage}[b]{.4\linewidth} % [b] => Ausrichtung an \caption
  
\begin{center}
\begin{tabular} {@{}l c @{}} \\ \hline
Nr. & Größe in cm\\
\hline
1 &160\\
2 &165\\
3 &165\\
4 & 165\\
5 & 173\\ 
6 & 174\\
7 & 176\\
8 & 180\\
9 & 180\\
10 &187\\
\hline
\end{tabular}
\vspace{1em}
\end{center}
  \end{minipage}
  %\hspace{.1\linewidth}% Abstand zwischen Bilder
  \begin{minipage}[b]{.4\linewidth} % [b] => Ausrichtung an \caption
  {{ \small
  \begin{align}
\bar{y} &= \dfrac{\sum y_i}{N} \nonumber \\
&= \frac{ 160 + 165 + 165 \dots + 187}{10} = 172,5 \nonumber
\end{align}
  }}
  \end{minipage}
\end{center}
}
\begin{frame}
\frametitle{Unterschiede der gewünschten Arbeitsstunden nach Geschlecht}
\begin{center}
\begin{tabular}{lrrr} \toprule 
Wert    & Männer & Frauen & Differenz \\ \midrule 
Modus   & 40     & 30 & 10\\
Median   & 40 & 30 & 10 \\
Arith. Mittel & 36.5 & 29.1 & 7.4 \\  \bottomrule
\end{tabular} 
\end{center}
\end{frame}
\begin{frame}
\frametitle{Gewünschte Arbeitsstunden nach Geschlecht über die Zeit}
\begin{center}
\resizebox{.7\linewidth}{!}{\includegraphics{lehre-gew-arbeitsstunden-mean-change-gender}}
\end{center}
\end{frame}

\section{Typische Unterschiedlichkeit}

\begin{frame}
  \frametitle{Typische Unterschiedlichkeit: Notation \& Begriffe}
\small
\begin{tabular}{lp{.15\linewidth}p{.3\linewidth}p{.35\linewidth}} \toprule 
Notation & Bezeichnung & Bedeutung & Beispiel\\ \midrule
$s^2$ & Varianz & Die durchschnittliche, quadrierte Abweichung der Werte zum Mittelwert. &  \\[.5em]
s & Standard\-abweichung & Die durchschnittliche Abweichung der Werte zum Mittelwert & Im Jahr 2018 unterscheiden sich die gewünschten Arbeitsstunden zwischen allen Befragten typischerweise um 10.4h. \\[.5em] 
D & Diversität & Die Wahrscheinlichkeit zwei Untersuchungseinheiten mit unterschiedlichen Ausprägungen zufällig aus der Stichprobe zu ziehen. & Die Wahrscheinlichkeit, im Jahr 2018 zwei zufällig gezogene Befragte mit unterschiedlichen Arbeitswünschen zu ziehen, beträgt  0.4.\\
 \bottomrule
\end{tabular}
  
\end{frame}
\frame{
\frametitle{Varianz}

Ein Maß für die Unterschiedlichkeit von Werten ist die \textit{Varianz}. Sie ist die \textit{durchschnittliche} quadrierte Abweichung zum arithmetischen Mittel (und damit eine Sonderform des arithmetischen Mittels). 

\begin{equation}
s^2 = \frac{\sum(y_i - \overline{y})^2}{n}
\end{equation}
}
\frame{
\frametitle{Varianz}

\begin{center}
  \begin{minipage}[b]{.2\linewidth} % [b] => Ausrichtung an \caption
  
\begin{center}
\scalebox{.7}{
\begin{tabular} {@{}l c @{}} \\ \hline
Nr. & Größe in cm\\
\hline
1 &160\\
2 &165\\
3 &165\\
4 & 165\\
5 & 173\\ 
6 & 174\\
7 & 176\\
8 & 180\\
9 & 180\\
10 &187\\
\hline
\end{tabular}}
\vspace{1em}
\end{center}
  \end{minipage}
  %\hspace{.1\linewidth}% Abstand zwischen Bilder
  \begin{minipage}[b]{.6\linewidth} % [b] => Ausrichtung an \caption
  {{ 
  \begin{align} 
s^2 &= \frac{\sum\limits_{i=1}^n(y_i - \overline{y})^2}{n} \nonumber \\
    &= \frac{(160 - 172,5)^2 + (165 - 172,5)^2 \dots + (187 - 172,5)^2}{10} = 66.25 \nonumber
\end{align}
  }}
  \end{minipage}
\end{center}
}
\frame{
\frametitle{Standardabweichung}

Wir wollen nicht wissen, was die durchschnittliche "`quadrierte Differenz"' ist, sondern, was die "`typischen Unterschiede"' der Befragten sind. Um eine sinnvolle, inhaltliche Interpretation zu ermöglichen, muss die Wurzel aus der Varianz gezogen werden. Diese Wurzel bezeichnet man als \textit{Standardabweichung}: 

\begin{equation}
  s = \sqrt{s^2}
\end{equation}

Da $\sqrt{66,25} \approx 8,14$ ist die typische Abweichung der Körpergrößen zwischen den beobachteten Personen 8,14 cm. 

}
\begin{frame}
\frametitle{Verteilung der Löhne in Deutschland 2018}
\begin{center}
\resizebox{.68\linewidth}{!}{
  \includegraphics{lohnvert2018-lehre}
}
\end{center}
\end{frame}

\subsection{Diversität}

\frame{
\frametitle{Diversität}

In einigen Fällen\footnote{Das sind vor allem: Verteilungen mit nominalen oder ordinal gemessenen Ausprägungen mit sehr unterschiedlichen Intervallen.} ist eine Berechnung eines Mittelwerts nicht sinnvoll. Dennoch möchten wir dann eine Möglichkeit haben, die Unterschiedlichkeit der verteilten Werte auszudrücken. Eine Möglichkeit, diese Unterschiedlichkeit zu messen, besteht in Diversitäts-Maßen, wie dem Simpson Diversitäts Index:  

\[ D = 1 - \sum\limits_{y}^Y P(y)^2 \]

y ist eine Ausprägung des Merkmals Y, Y die Anzahl aller Ausprägungen, und N die Gesamtanzahl aller Beobachtungen ist. Der Index drückt die Wahrscheinlichkeit aus, durch eine Ziehung zwei \textit{unterschiedlich ausgeprägte} Beobachtungen zu erhalten. D schwankt zwischen 0 und 1. Ist D = 1, liegt höchste Diversität vor. 

}
\begin{frame}
\frametitle{Belastung durch Arbeitsstunden}
\small
\begin{columns}
\begin{column}{.4\linewidth}
\begin{tabular}{lrrr} \toprule
Belastung & Männer & Frauen & Gesamt\\ \midrule
Gar nicht &       854  &       747 &     1,601 \\
Mäßig &     4,033  &     3,399 &     7,432 \\
Stark &     2,459  &     2,186 &     4,645 \\
Sehr stark&       588 &        553 &     1,141 \\ \midrule
Gesamt     & 7,934    &  6,885 &    14,819 \\ \bottomrule
\end{tabular}
\end{column}

\begin{column}{.5\linewidth}
Die Tabelle stellt Angabe der subjektiven Belastung durch die Arbeitsstunden ab. Die Diversität der Angaben über alle Befragten lässt sich wie folgt berechnen:\\[1em]

\vfill 

\resizebox{1\linewidth}{!}{
$ D = 1 - \left( (\dfrac{1601}{14819})^2 + (\dfrac{7432}{14819})^2 + (\dfrac{4645}{14819})^2 + (\dfrac{1141}{14819})^2 \right) = 0.633 $
}
\quad \newline
\vfill

Würden wir zwei Befragte zufällig ziehen, erhielten wir daher mit einer Wahrscheinlichkeit von 63.3\% unterschiedliche Angaben. Die Diversität innerhalb von Männern ist 0.628 und innerhalb von Frauen 0.637. 

\end{column}

\end{columns}

\end{frame}
\frame{
\frametitle{Ethnische Diversität}
\begin{center}
\resizebox{.7\linewidth}{!}{\includegraphics{US-race-diversity}}
\end{center}

}